\documentclass[11pt]{article}

\usepackage[margin=1in]{geometry}
\usepackage{array}

\title{Investigacion Ley de Conway}
\author{Andreas Garoufalias}
\date{\today}

\begin{document}

\maketitle

\begin{abstract}
	Este informe explica la ley de conway.
\end{abstract}

\section{Ley de Conway}

\textit{Any organisation that designs a system (defined broadly) will produce a design whose structure is a copy of the organisation's communication structure. -- Melvin Conway}

Cualquier organizacion que disena un sistema (definida ampliamente) va a producir un diseno cuya estructura es una copia de la estructura de comunicacion de la organizacion.

\section{Explicacion}

La ley de conway, en esencia, es la observacion que las arquitecturas de los sistemas de software son parecidas a la organizacion de los equipos de desarrollo que lo construyen.

Aunque este concepto se habla usualmente con respecto a software, se puede aplicar a sistemas en general.

\section{Aspectos}

\begin{itemize}
	\item Refleccion de Estructura
		\begin{itemize}
			\item Si una organizacion esta dividida en silos, el sistema resultante tiene una probabilidad alta de terminar siendo construida en silos con vinculos entre equipos.
		\end{itemize}
		
	\item Cuellas de botello de comunicacion
		\begin{itemize}
			\item Comunicacion mala entre equipos resulta en malas implementaciones en sistemas.
		\end{itemize}

	\item Maneobra inversa de Conway
		\begin{itemize}
			\item Organizaciones pueden estructurar sus equipos para igualar la arquitectura del sistema de software que quieren.
		\end{itemize}
	
	\item Dimensiones del equipo
		\begin{itemize}
			\item Equipos diminutos tienden a crear aplicaciones mas simples, cohesivos, o monoliticos, mientras que equipos grandes tienden a producir sistemas mas fragmentados o complejos.
		\end{itemize}
\end{itemize}

\end{document}